\documentclass[10pt,a4paper]{amsart}
\usepackage{amsmath,amssymb,amsthm,graphicx}
\usepackage{verbatim}
\newenvironment{reference}{\comment}{\endcomment}
%\newenvironment{reference}{}{}
\newenvironment{slogan}{\comment}{\endcomment}
\newenvironment{history}{\comment}{\endcomment}

% For commutative diagrams we use Xy-pic
\usepackage[all]{xy}

% We use 2cell for 2-commutative diagrams.
\xyoption{2cell}
% Theorem environments.
%
\theoremstyle{plain}
\newtheorem{theorem}[subsection]{Theorem}
\newtheorem{proposition}[subsection]{Proposition}
\newtheorem{lemma}[subsection]{Lemma}

\theoremstyle{definition}
\newtheorem{definition}[subsection]{Definition}
\newtheorem{example}[subsection]{Example}
\newtheorem{exercise}[subsection]{Exercise}
\newtheorem{situation}[subsection]{Situation}

\theoremstyle{remark}
\newtheorem{remark}[subsection]{Remark}
\newtheorem{remarks}[subsection]{Remarks}

\numberwithin{equation}{subsection}

% Macros
%
\def\lim{\mathop{\mathrm{lim}}\nolimits}
\def\colim{\mathop{\mathrm{colim}}\nolimits}
\def\Spec{\mathop{\mathrm{Spec}}}
\def\Hom{\mathop{\mathrm{Hom}}\nolimits}
\def\Ext{\mathop{\mathrm{Ext}}\nolimits}
\def\SheafHom{\mathop{\mathcal{H}\!\mathit{om}}\nolimits}
\def\SheafExt{\mathop{\mathcal{E}\!\mathit{xt}}\nolimits}
\def\Sch{\mathit{Sch}}
\def\Mor{\mathop{\mathrm{Mor}}\nolimits}
\def\Ob{\mathop{\mathrm{Ob}}\nolimits}
\def\Sh{\mathop{\mathit{Sh}}\nolimits}
\def\NL{\mathop{N\!L}\nolimits}
\def\CH{\mathop{\mathrm{CH}}\nolimits}
\def\proetale{{pro\text{-}\acute{e}tale}}
\def\etale{{\acute{e}tale}}
\def\QCoh{\mathit{QCoh}}
\def\Ker{\mathop{\mathrm{Ker}}}
\def\Im{\mathop{\mathrm{Im}}}
\def\Coker{\mathop{\mathrm{Coker}}}
\def\Coim{\mathop{\mathrm{Coim}}}

% Boxtimes
%
\DeclareMathSymbol{\boxtimes}{\mathbin}{AMSa}{"02}

%
% Macros for moduli stacks/spaces
%
\def\QCohstack{\mathcal{QC}\!\mathit{oh}}
\def\Cohstack{\mathcal{C}\!\mathit{oh}}
\def\Spacesstack{\mathcal{S}\!\mathit{paces}}
\def\Quotfunctor{\mathrm{Quot}}
\def\Hilbfunctor{\mathrm{Hilb}}
\def\Curvesstack{\mathcal{C}\!\mathit{urves}}
\def\Polarizedstack{\mathcal{P}\!\mathit{olarized}}
\def\Complexesstack{\mathcal{C}\!\mathit{omplexes}}
% \Pic is the operator that assigns to X its picard group, usage \Pic(X)
% \Picardstack_{X/B} denotes the Picard stack of X over B
% \Picardfunctor_{X/B} denotes the Picard functor of X over B
\def\Pic{\mathop{\mathrm{Pic}}\nolimits}
\def\Picardstack{\mathcal{P}\!\mathit{ic}}
\def\Picardfunctor{\mathrm{Pic}}
\def\Deformationcategory{\mathcal{D}\!\mathit{ef}}
\usepackage{cleveref}
\crefname{originallemma}{original lemma}{original lemmas}
\Crefname{originallemma}{Original lemma}{Original lemmas}
\makeatletter
\newlabel{algebra-lemma-derivative-zero-pth-power}{{031W}{source}{}{sourcepointer-1}{}}
\newlabel{algebra-lemma-derivative-zero-pth-power@cref}{{[originallemma][0][]031W}{[1][0][]source}{}{}{}}
\makeatother
\crefname{originalreference}{original reference}{original references}
\Crefname{originalreference}{Original reference}{Original references}
\makeatletter
\newlabel{lemma-bound-dualizing-implies-bound}{{0B4S}{source}{}{sourcepointer-2}{}}
\newlabel{lemma-bound-dualizing-implies-bound@cref}{{[originalreference][0][]0B4S}{[1][0][]source}{}{}{}}
\makeatother
\crefname{originalreference}{original reference}{original references}
\Crefname{originalreference}{Original reference}{Original references}
\makeatletter
\newlabel{lemma-compare-differentials-dualizing}{{0B4T}{source}{}{sourcepointer-3}{}}
\newlabel{lemma-compare-differentials-dualizing@cref}{{[originalreference][0][]0B4T}{[1][0][]source}{}{}{}}
\makeatother
\crefname{originalreference}{original reference}{original references}
\Crefname{originalreference}{Original reference}{Original references}
\makeatletter
\newlabel{lemma-dominate-by-blowing-up-in-points}{{0AHI}{source}{}{sourcepointer-4}{}}
\newlabel{lemma-dominate-by-blowing-up-in-points@cref}{{[originalreference][0][]0AHI}{[1][0][]source}{}{}{}}
\makeatother
\crefname{originalreference}{original reference}{original references}
\Crefname{originalreference}{Original reference}{Original references}
\makeatletter
\newlabel{lemma-go-up-separable}{{0AXL}{source}{}{sourcepointer-5}{}}
\newlabel{lemma-go-up-separable@cref}{{[originalreference][0][]0AXL}{[1][0][]source}{}{}{}}
\makeatother
\crefname{originalreference}{original reference}{original references}
\Crefname{originalreference}{Original reference}{Original references}
\makeatletter
\newlabel{lemma-trace-extends}{{0AX6}{source}{}{sourcepointer-6}{}}
\newlabel{lemma-trace-extends@cref}{{[originalreference][0][]0AX6}{[1][0][]source}{}{}{}}
\makeatother
\crefname{originalreference}{original reference}{original references}
\Crefname{originalreference}{Original reference}{Original references}
\makeatletter
\newlabel{lemma-trace-higher}{{0AX5}{source}{}{sourcepointer-7}{}}
\newlabel{lemma-trace-higher@cref}{{[originalreference][0][]0AX5}{[1][0][]source}{}{}{}}
\makeatother
\crefname{originallemma}{original lemma}{original lemmas}
\Crefname{originallemma}{Original lemma}{Original lemmas}
\makeatletter
\newlabel{more-algebra-lemma-power-series-ring-subfields}{{07P5}{source}{}{sourcepointer-8}{}}
\newlabel{more-algebra-lemma-power-series-ring-subfields@cref}{{[originallemma][0][]07P5}{[1][0][]source}{}{}{}}
\makeatother
\title{Reduction to rational singularities for the integral closure in a purely inseparable extension of degree p}
\author{The Stacks Project Authors}
\date{stacks-project:resolve.tex:local-snapshot-20261012; GNU FDL 1.2 or later}
\begin{document}
\maketitle
\noindent\textit{Editorial note.} All delivered mathematical text is the author's own, copied from the pinned Stacks Project chapter file resolve.tex; no statement or proof step was rewritten, reordered or removed.\par\medskip
\noindent\textit{Editorial note.} The counted proof is the author's complete argument: the choice of the finite subextension of the coefficient field with the computation of the rank of the relative differential module, the construction of the comparison map from the trace together with the explicit computation on the power basis showing that its cokernel is annihilated by a nonzero element, and the transport of the comparison along modifications reduced to sequences of point blowups followed by the bound criterion.\par\medskip
\noindent\textit{Editorial note.} The label lemma-go-up-degree-p names the counted statement, and the proof environment immediately following it in the source is delivered verbatim as its proof.\par\medskip
\noindent\textit{Editorial note.} The counted statement is carried by the author's own lemma environment in the chapter source.\par\medskip
\noindent\textit{Editorial note.} Pre-open gate: python operations/agent-authoring/tools/precard\_gate.py --source resolve.tex --label lemma-go-up-degree-p --start 2348 --end 2439 --local-id 28071 returned RESULT=OK (L28071 free, occupancy FREE, public interval clear). Pre-assembly tag self-check: canonical tag 0B4U, not present in reference/public-items.jsonl.\par\medskip
\noindent\textit{Editorial note.} The statement and proof are copied from the pinned chapter file; the cited inputs are kept as external pointers with their canonical Stacks tags.\par\medskip
\noindent\textit{Allowed standard prerequisites (editorial).}\begin{itemize}
\item Description of the trace map on differential forms for a finite extension. Same chapter, the lemma extending the trace to differential forms and the lemma giving its explicit description on top forms. Supplies the trace maps used to build the comparison map and the displayed computation on the basis elements.
\item Finiteness of an inseparable extension of the coefficient field inside the power series ring. More on Algebra, the lemma producing a finite subextension of the coefficient field over which the differential of the defining element remains nonzero. Supplies the field k' and hence the regular local subring A\_0 over which the relative differential module is finite free of known rank.
\item Reduction of a normal modification to a finite one over a sequence of point blowups. Same chapter, the corresponding reduction for the separable case and the lemma dominating a modification by a sequence of blowups in closed points. Supplies the commutative square in which the lower horizontal map is a sequence of blowups in closed points.
\item Comparison of the dualizing modules along a sequence of point blowups and the bound criterion. Same chapter, the lemma comparing differentials with the dualizing module for such a sequence, and the lemma concluding the bound on the dualizing module from the existence of such a comparison. Supply the second arrow of the composite and the final step deducing the bound on the dualizing module of B.
\end{itemize}\medskip
\ifdefined\hypertarget\hypertarget{sourcepointer-1}{}\fi
\noindent\textit{Source reference.} Original lemma 031W: sources/stacks-chow/tags.txt, source lines 3999--3999. Prerequisite pointer invoked at proof line(s) 2359; it names a standard fact of a sibling chapter that the counted argument uses, and every step of the counted proof built on it is retained verbatim in the delivered body.\par\medskip
\ifdefined\hypertarget\hypertarget{sourcepointer-2}{}\fi
\noindent\textit{Source reference.} Original reference 0B4S: sources/stacks-resolve/pinned-original-source.tex, source lines 2233--2233. Prerequisite pointer invoked at proof line(s) 2438; it names a standard fact of the same chapter that the counted argument uses, and every step of the counted proof built on it is retained verbatim in the delivered body.\par\medskip
\ifdefined\hypertarget\hypertarget{sourcepointer-3}{}\fi
\noindent\textit{Source reference.} Original reference 0B4T: sources/stacks-resolve/pinned-original-source.tex, source lines 2291--2291. Prerequisite pointer invoked at proof line(s) 2428; it names a standard fact of the same chapter that the counted argument uses, and every step of the counted proof built on it is retained verbatim in the delivered body.\par\medskip
\ifdefined\hypertarget\hypertarget{sourcepointer-4}{}\fi
\noindent\textit{Source reference.} Original reference 0AHI: sources/stacks-resolve/pinned-original-source.tex, source lines 1146--1146. Prerequisite pointer invoked at proof line(s) 2409; it names a standard fact of the same chapter that the counted argument uses, and every step of the counted proof built on it is retained verbatim in the delivered body.\par\medskip
\ifdefined\hypertarget\hypertarget{sourcepointer-5}{}\fi
\noindent\textit{Source reference.} Original reference 0AXL: sources/stacks-resolve/pinned-original-source.tex, source lines 2139--2139. Prerequisite pointer invoked at proof line(s) 2407; it names a standard fact of the same chapter that the counted argument uses, and every step of the counted proof built on it is retained verbatim in the delivered body.\par\medskip
\ifdefined\hypertarget\hypertarget{sourcepointer-6}{}\fi
\noindent\textit{Source reference.} Original reference 0AX6: sources/stacks-resolve/pinned-original-source.tex, source lines 591--591. Prerequisite pointer invoked at proof line(s) 2378; it names a standard fact of the same chapter that the counted argument uses, and every step of the counted proof built on it is retained verbatim in the delivered body.\par\medskip
\ifdefined\hypertarget\hypertarget{sourcepointer-7}{}\fi
\noindent\textit{Source reference.} Original reference 0AX5: sources/stacks-resolve/pinned-original-source.tex, source lines 501--501. Prerequisite pointer invoked at proof line(s) 2403; it names a standard fact of the same chapter that the counted argument uses, and every step of the counted proof built on it is retained verbatim in the delivered body.\par\medskip
\ifdefined\hypertarget\hypertarget{sourcepointer-8}{}\fi
\noindent\textit{Source reference.} Original lemma 07P5: sources/stacks-chow/tags.txt, source lines 9785--9785. Prerequisite pointer invoked at proof line(s) 2361; it names a standard fact of a sibling chapter that the counted argument uses, and every step of the counted proof built on it is retained verbatim in the delivered body.\par\medskip
\begin{lemma}
\label{lemma-go-up-degree-p}
Let $p$ be a prime number. Let $A$ be a complete regular local ring of
dimension $2$ and characteristic $p$. Let $L/K$ be a degree $p$ inseparable
extension of the fraction field $K$ of $A$. Let $B \subset L$ be the integral
closure of $A$. Then reduction to rational singularities is possible for $B$.
\end{lemma}

\begin{proof}
We have $A = k[[x, y]]$. Write $L = K[x]/(x^p - f)$ for some $f \in A$
and denote $g \in B$ the congruence class of $x$, i.e., the element such
that $g^p = f$. By Algebra, Lemma \ref{algebra-lemma-derivative-zero-pth-power}
we see that $\text{d}f$ is nonzero in $\Omega_{K/\mathbf{F}_p}$. By
More on Algebra, Lemma \ref{more-algebra-lemma-power-series-ring-subfields}
there exists a subfield $k^p \subset k' \subset k$ with
$p^e = [k : k'] < \infty$ such that $\text{d}f$ is nonzero
in $\Omega_{K/K_0}$ where $K_0$ is the fraction field of
$A_0 = k'[[x^p, y^p]] \subset A$. Then
$$
\Omega_{A/A_0} =
A \otimes_k \Omega_{k/k'} \oplus A \text{d}x \oplus A \text{d}y
$$
is finite free of rank $e + 2$. Set $\omega_A = \Omega^{e + 2}_{A/A_0}$.
Consider the canonical map
$$
\text{Tr} :
\Omega^{e + 2}_{B/A_0}
\longrightarrow
\Omega^{e + 2}_{A/A_0} = \omega_A
$$
of Lemma \ref{lemma-trace-extends}. By duality this determines a map
$$
c : \Omega^{e + 2}_{B/A_0} \to \omega_B = \Hom_A(B, \omega_A)
$$
Claim: the cokernel of $c$ is annihilated by a nonzero element of $B$.

\medskip\noindent
Since $\text{d}f$ is nonzero in $\Omega_{A/A_0}$ we can find
$\eta_1, \ldots, \eta_{e + 1} \in \Omega_{A/A_0}$ such that
$\theta = \eta_1 \wedge \ldots \wedge \eta_{e + 1} \wedge \text{d}f$ is
nonzero in $\omega_A = \Omega^{e + 2}_{A/A_0}$. To prove the claim we
will construct elements $\omega_i$ of $\Omega^{e + 2}_{B/A_0}$,
$i = 0, \ldots, p - 1$ which are mapped to
$\varphi_i \in \omega_B = \Hom_A(B, \omega_A)$
with $\varphi_i(g^j) = \delta_{ij}\theta$ for $j = 0, \ldots, p - 1$.
Since $\{1, g, \ldots, g^{p - 1}\}$ is a basis for $L/K$ this
proves the claim. We set
$\eta = \eta_1 \wedge \ldots \wedge \eta_{e + 1}$
so that $\theta = \eta \wedge \text{d}f$.
Set $\omega_i = \eta \wedge g^{p - 1 - i}\text{d}g$. Then
by construction we have
$$
\varphi_i(g^j) = \text{Tr}(g^j \eta \wedge g^{p - 1 - i}\text{d}g) =
\text{Tr}(\eta \wedge g^{p - 1 - i + j}\text{d}g) = \delta_{ij} \theta
$$
by the explicit description of the trace map in Lemma \ref{lemma-trace-higher}.

\medskip\noindent
Let $Y \to \Spec(B)$ be a normal modification. Exactly as in the proof of
Lemma \ref{lemma-go-up-separable} we can reduce to the case where $Y$
is finite over a modification $X$ of $\Spec(A)$. By
Lemma \ref{lemma-dominate-by-blowing-up-in-points} we may
even assume $X \to \Spec(A)$ is the result of a sequence of
blowing ups in closed points. Picture:
$$
\xymatrix{
Y \ar[d]_\pi \ar[r]_-g & \Spec(B) \ar[d] \\
X \ar[r]^-f & \Spec(A)
}
$$
We may apply Lemma \ref{lemma-trace-extends} to $\pi$ and
we obtain the first arrow in
$$
\pi_*(\Omega^{e + 2}_{Y/A_0})
\xrightarrow{\text{Tr}}
(\Omega^{e + 2}_{X/A_0})^{**}
\xrightarrow{\varphi_X}
\omega_X
$$
and the second arrow is from
Lemma \ref{lemma-compare-differentials-dualizing}
(because $f$ is a sequence of blowups in closed points).
By duality for the finite morphism $\pi$ this corresponds to a map
$$
c_Y : \Omega^{e + 2}_{Y/A_0} \longrightarrow \omega_Y
$$
extending the map $c$ above. Hence we see that the image of
$\Gamma(Y, \omega_Y) \to \omega_B$ contains the image of $c$.
By our claim we see that the cokernel is annihilated by
a fixed nonzero element of $B$. We conclude by
Lemma \ref{lemma-bound-dualizing-implies-bound}.
\end{proof}

\end{document}
